\subsection{Characteristic and Minimal Polynomial Foundations}

\begin{motivation} \label{mot:characteristic-and-minimal-polynomials} \hlt{Two Spectral Polynomials}\\
The characteristic polynomial records the whole primary decomposition with dimensions.\\
The minimal polynomial records the shortest polynomial relation satisfied by the operator.
\end{motivation}

\begin{definition} \label{def:matrix-characteristic-polynomial} \hlt{Characteristic Polynomial}\\
For $A\in\operatorname{Mat}(n,n,\mathbb{F})$, the \hlt{characteristic polynomial} of $A$ is $p_A(z)\defeq\det(zI_n-A)$.
\end{definition}

\begin{theorem} \label{thm:matrix-characteristic-polynomial-basic-coefficients} \hlt{Basic Form of the Characteristic Polynomial}\\
Let $A=[a_{ij}]\in\operatorname{Mat}(n,n,\mathbb{F})$. Then $p_A$ is monic of degree $n$ and $p_A(z)=z^n-(\operatorname{tr}A)z^{n-1}+\cdots+(-1)^n\det A$.\\
Moreover, for $\lambda\in\mathbb{F}$, $\lambda\in\operatorname{spec}(A)$ if and only if $p_A(\lambda)=0$.
\begin{proof}
The only determinant term of degree greater than $n-2$ comes from the identity permutation:
\begin{align}
(z-a_{11})(z-a_{22})\cdots(z-a_{nn})=z^n-(a_{11}+\cdots+a_{nn})z^{n-1}+\cdots.
\nonumber
\end{align}
Thus $p_A$ is monic and the coefficient of $z^{n-1}$ is $-\operatorname{tr}A$; also $p_A(0)=\det(-A)=(-1)^n\det A$. Finally, $p_A(\lambda)=0$ if and only if $\lambda I_n-A$ is singular, equivalently $\ker(A-\lambda I_n)\neq\{\mathbf{0}\}$.\\
\hlt{Structural Commentary}: The diagonal product determines the top two coefficients.\\
\hlt{Key Insight}: Eigenvalues are precisely the roots of $\det(zI_n-A)$.
\end{proof}
\end{theorem}

\begin{definition} \label{def:matrix-algebraic-multiplicity} \hlt{Algebraic Multiplicity}\\
For $\lambda\in\mathbb{F}$, let $m_A(\lambda)$ be the multiplicity of $\lambda$ as a root of $p_A$, with $m_A(\lambda)=0$ if $\lambda\notin\operatorname{spec}(A)$.\\
For $\lambda\in\operatorname{spec}(A)$, this is the \hlt{algebraic multiplicity} of $\lambda$. Unqualified multiplicity means algebraic multiplicity.
\end{definition}

\begin{figure}[H]
  \centering
  \begin{tikzpicture}[
  >=Stealth,
  node distance=0.85cm and 1.0cm,
  box/.style={draw=blue!55!black, fill=blue!5, rounded corners=2pt, align=center, minimum height=0.78cm, minimum width=2.05cm, font=\small},
  root/.style={circle, draw=red!70!black, fill=red!10, inner sep=1.5pt, font=\small},
  arr/.style={->, line width=0.65pt, draw=black!70}
]
  \node[box] (matrix) {$A$};
  \node[box, right=of matrix] (poly) {$p_A$};
  \node[box, right=of poly] (roots) {roots};

  \draw[arr] (matrix) -- (poly);
  \draw[arr] (poly) -- (roots);

  \node[right=0.72cm of roots, inner sep=0pt, font=\small] (openroots) {$\{$};
  \node[root, right=0.10cm of openroots] (l1) {$\lambda_1$};
  \node[root, right=0.24cm of l1] (l2) {$\lambda_2$};
  \node[root, right=0.24cm of l2] (ld) {$\lambda_d$};
  \node[right=0.10cm of ld, inner sep=0pt, font=\small] (closeroots) {$\}$};
  \draw[arr] (roots) -- (openroots);
\end{tikzpicture}

  \caption{Form $p_A(z)=\det(zI_n-A)$ and factor it; each root $\lambda$ is counted by $m_A(\lambda)$.}
  \label{fig:characteristic-polynomial-roots}
\end{figure}

\begin{method} \label{meth:characteristic-equation-eigenspace-computation} \hlt{Characteristic Equation and Eigenspace Computation}\\
To compute eigenvalues and eigenspaces of $A\in\operatorname{Mat}(n,n,\mathbb{F})$:
\begin{enumerate}[label=(\roman*)]
\setlength{\itemsep}{0pt}
\item Form the characteristic polynomial $p_A(z)=\det(zI_n-A)$ from Definition \ref{def:matrix-characteristic-polynomial}.
\item Solve the characteristic equation $p_A(\lambda)=0$. By Theorem \ref{thm:matrix-characteristic-polynomial-basic-coefficients}, the solutions are exactly the eigenvalues in $\mathbb{F}$.
\item For each eigenvalue $\lambda$, row-reduce $A-\lambda I_n$ and compute
\begin{align}
E_{\lambda}(A)=\ker(A-\lambda I_n). \nonumber
\end{align}
\item Choose a basis of each eigenspace. The dimension of this basis is the geometric multiplicity of $\lambda$.
\end{enumerate}
Algebraic multiplicity comes from the factorisation of $p_A$, while geometric multiplicity comes from the nullspace computation.
\end{method}

\begin{proposition} \label{prop:upper-triangular-characteristic-polynomial} \hlt{Upper-Triangular Eigenvalues}\\
If $A=[a_{ij}]\in\operatorname{Mat}(n,n,\mathbb{F})$ is upper triangular, then $p_A(z)=\prod\limits_{j=1}^{n}(z-a_{jj})$.\\
Thus $\operatorname{spec}(A)=\{a_{jj}:1\leq j\leq n\}$, and for $\lambda\in\operatorname{spec}(A)$, $m_A(\lambda)=\#\{j:a_{jj}=\lambda\}$.
\begin{proof}
Since $zI_n-A$ is upper triangular with diagonal entries $z-a_{jj}$, $p_A(z)=\det(zI_n-A)=\prod\limits_{j=1}^{n}(z-a_{jj})$.\\
Thus the roots of $p_A$ are exactly the diagonal entries, and the exponent of $(z-\lambda)$ is $\#\{j:a_{jj}=\lambda\}$.\\
\hlt{Structural Commentary}: Triangular form removes all off-diagonal determinant contributions.\\
\hlt{Key Insight}: Once a matrix is triangular, the spectrum is visible on the diagonal.
\end{proof}
\end{proposition}

\begin{theorem} \label{thm:similarity-preserves-characteristic-polynomial} \hlt{Similarity Preserves Characteristic Data}\\
Let $A,B\in\operatorname{Mat}(n,n,\mathbb{F})$ and suppose $B=SAS^{-1}$ for some invertible $S$. Then $p_B=p_A$. In particular, for each $\lambda\in\operatorname{spec}(A)$, $m_B(\lambda)=m_A(\lambda)$ and $\dim E_{\lambda}(B)=\dim E_{\lambda}(A)$.
\begin{proof}
For every $z$, $p_B(z) =\det(zI_n-SAS^{-1}) =\det(S(zI_n-A)S^{-1}) =\det(S)\det(zI_n-A)\det(S)^{-1} =p_A(z)$.\\
Thus roots and root multiplicities agree, so $m_B(\lambda)=m_A(\lambda)$. Also $B-\lambda I_n=S(A-\lambda I_n)S^{-1}$.\\
Hence $\mathbf{x}\in\ker(A-\lambda I_n)$ if and only if $S\mathbf{x}\in\ker(B-\lambda I_n)$, $S$ restricts to isomorphism between two eigenspaces.\\
\hlt{Structural Commentary}: Similarity changes coordinates, not the shifted determinant or shifted kernel dimension.\\
\hlt{Key Insight}: Conjugation inserts $S$ and $S^{-1}$, which cancel in determinants and identify kernels.
\end{proof}
\end{theorem}

\begin{proposition} \label{prop:block-upper-triangular-characteristic-polynomial} \hlt{Block Upper-Triangular Characteristic Polynomial}\\
Let $A=
\begin{bmatrix}
B & C\\
0 & D
\end{bmatrix}$, $B$ and $D$ are square. Then $p_A(z)=p_B(z)p_D(z)$, so $m_A(\lambda)=m_B(\lambda)+m_D(\lambda) \ \forall\lambda\in\mathbb{F}$.
\begin{proof}
The shifted matrix is block upper triangular, $zI-A=
\begin{bmatrix}
zI-B & -C\\
0 & zI-D \end{bmatrix}$.\\
Theorem \ref{thm:block-upper-triangular-determinant} gives $p_A(z)=\det(zI-B)\det(zI-D)=p_B(z)p_D(z)$.\\
Roots of a product are the roots of the factors, with multiplicity added.\\
\hlt{Structural Commentary}: The off-diagonal block couples vectors but not determinant factors.\\
\hlt{Key Insight}: In block upper-triangular form, spectral data lives on the diagonal blocks.
\end{proof}
\end{proposition}

\begin{theorem} \label{thm:geometric-multiplicity-bounded-by-algebraic-multiplicity} \hlt{Geometric Multiplicity Is at Most Algebraic Multiplicity}\\
Let $A\in\operatorname{Mat}(n,n,\mathbb{F})$ and let $\lambda\in\operatorname{spec}(A)$. Then $\dim E_{\lambda}(A)\leq m_A(\lambda)$.
\begin{proof}
Let $k=\dim E_{\lambda}(A)$ and choose a basis $\mathbf{x}_1,\ldots,\mathbf{x}_k$ of $E_{\lambda}(A)$. Extend it to a basis of $\mathbb{F}^n$, and let $S$ be the matrix whose columns are this basis. In the new basis,
\begin{align}
S^{-1}AS=
\begin{bmatrix}
\lambda I_k & *\\
0 & C
\end{bmatrix}.
\nonumber
\end{align}
By Theorem \ref{thm:similarity-preserves-characteristic-polynomial}, Proposition \ref{prop:block-upper-triangular-characteristic-polynomial}, $p_A(z)=p_{S^{-1}AS}(z)=(z-\lambda)^kp_C(z)$. Thus $m_A(\lambda)\geq k=\dim E_{\lambda}(A)$.\\
\hlt{Structural Commentary}: An eigenspace-first basis exposes a $\lambda I_k$ block.\\
\hlt{Key Insight}: Each independent eigenvector forces one visible factor $z-\lambda$ in the characteristic polynomial.
\end{proof}
\end{theorem}

\begin{theorem} \label{thm:diagonalizable-iff-algebraic-geometric-multiplicities-match} \hlt{Diagonalizability and Multiplicity}\\
Let $A\in\operatorname{Mat}(n,n,\mathbb{C})$. Then $A$ is diagonalizable if and only if $\dim E_{\lambda}(A)=m_A(\lambda)$ for every $\lambda\in\operatorname{spec}(A)$.
\begin{proof}
\begin{enumerate}[label=\roman*.]
\setlength{\itemsep}{0pt}
\item[$(\Rightarrow)$] If $A$ is diagonalizable, then $A$ is similar to a diagonal matrix $D$. Theorem \ref{thm:similarity-preserves-characteristic-polynomial} gives $m_A(\lambda)=m_D(\lambda)$ and $\dim E_{\lambda}(A)=\dim E_{\lambda}(D)$. For diagonal $D$, both numbers count diagonal entries equal to $\lambda$.
\item[$(\Leftarrow)$] Assume $\dim E_{\lambda}(A)=m_A(\lambda)$ for every eigenvalue.\\
If $\mu_1,\ldots,\mu_d$ are distinct eigenvalues, then $\sum\limits_{i=1}^{d}m_A(\mu_i)=n$. Hence $\dim E_{\mu_1}(A)+\cdots+\dim E_{\mu_d}(A)=n$.\\
Theorem \ref{thm:conditions-equivalent-to-diagonalizability} gives an eigenbasis, so $A$ is diagonalizable.
\end{enumerate}
\hlt{Structural Commentary}: Diagonalizability means root counts match eigenspace dimensions.\\
\hlt{Key Insight}: No missing eigendirections means an eigenbasis.
\end{proof}
\end{theorem}

\begin{definition} \label{def:operator-characteristic-polynomial} \hlt{Characteristic Polynomial of an Operator}\\
Let $\mathcal{V}\neq\{\mathbf{0}\}$ be finite-dimensional over $\mathbb{C}$, and let $T\in\LL(\mathcal{V})$. If $\lambda_1,\ldots,\lambda_m$ are the distinct eigenvalues of $T$ and $d_j=\dim \mathcal{G}_{\lambda_j}(T)$, then the \hlt{characteristic polynomial} of $T$ is $q_T(z)\defeq\prod\limits_{j=1}^{m}(z-\lambda_j)^{d_j}$.
\end{definition}

\begin{proposition} \label{prop:operator-characteristic-polynomial-degree-roots} \hlt{Degree and Roots of $q_T$}\\
Let $\mathcal{V}\neq\{\mathbf{0}\}$ be finite-dimensional over $\mathbb{C}$, let $T\in\LL(\mathcal{V})$, and let $q_T$ be the characteristic polynomial of $T$.
\begin{enumerate}[label=(\roman*)]
\setlength{\itemsep}{0pt}
\item $\deg q_T=\dim(\mathcal{V})$.
\item The roots of $q_T$ are precisely the eigenvalues of $T$.
\end{enumerate}
\begin{proof}
Write $q_T(z)=\prod\limits_{j=1}^{m}(z-\lambda_j)^{d_j}$, where $\lambda_1,\ldots,\lambda_m$ are the distinct eigenvalues and $d_j=\dim \mathcal{G}_{\lambda_j}(T)$.
\begin{enumerate}[label=(\roman*)]
\setlength{\itemsep}{0pt}
\item Proposition \ref{prop:sum-of-complex-eigenvalue-multiplicities} gives $\deg q_T=\sum_jd_j=\dim(\mathcal{V})$.
\item The displayed factorization has exactly the roots $\lambda_1,\ldots,\lambda_m$.
\end{enumerate}
\hlt{Structural Commentary}: Each primary summand contributes one root factor.\\
\hlt{Key Insight}: Multiplicity is dimension.
\end{proof}
\end{proposition}

\begin{figure}[H]
  \centering
  \begin{tikzpicture}[
  >=Stealth,
  every node/.style={font=\small},
  block/.style={draw=blue!60!black, fill=blue!6, rounded corners=2pt, align=center, inner xsep=4.5pt, inner ysep=4pt, minimum width=1.7cm},
  factor/.style={draw=dartmouthgreen!70!black, fill=dartmouthgreen!8, rounded corners=2pt, align=center, inner xsep=4.5pt, inner ysep=4pt, minimum width=1.7cm},
  result/.style={draw=gray!60, fill=gray!6, rounded corners=2pt, align=center, inner xsep=5pt, inner ysep=4.5pt},
  arr/.style={->, line width=0.7pt, draw=black!70}
]
  \node[block] (G1) at (-2.55,0) {$G(\lambda_1,T)$};
  \node (Gdots) at (0,0) {$\cdots$};
  \node[block] (Gm) at (2.55,0) {$G(\lambda_m,T)$};

  \node[factor] (F1) at (-2.55,-1.08) {$(z-\lambda_1)^{d_1}$};
  \node (Fdots) at (0,-1.08) {$\cdots$};
  \node[factor] (Fm) at (2.55,-1.08) {$(z-\lambda_m)^{d_m}$};

  \node[result] (q) at (0,-2.10) {$q_T(z)=\prod_{j=1}^{m}(z-\lambda_j)^{d_j}$};

  \draw[arr] (G1) -- (F1);
  \draw[arr] (Gm) -- (Fm);
  \draw[arr] (F1.south east) -- (q.north west);
  \draw[arr] (Fm.south west) -- (q.north east);
\end{tikzpicture}

  \caption{Each primary summand contributes $(z-\lambda_j)^{d_j}$, where $d_j=\dim \mathcal{G}_{\lambda_j}(T)$.}
  \label{fig:operator-characteristic-polynomial-primary-blocks}
\end{figure}

\begin{definition} \label{def:monic-polynomial} \hlt{Monic Polynomial}\\
A nonzero polynomial is \hlt{monic} if its highest-degree coefficient is $1$.
\end{definition}

\begin{theorem} \label{thm:operator-minimal-polynomial-existence} \hlt{Existence of the Minimal Polynomial}\\
Let $\mathcal{V}\neq\{\mathbf{0}\}$ be finite-dim over $\mathbb{F}$ and let $T\in\LL(\mathcal{V})$. Then $T$ has a unique least-degree monic annihilator $\mu_T$.
\begin{proof}
Proposition \ref{prop:finite-dimensional-operators-have-annihilating-polynomials} gives a nonzero annihilator; scaling makes it monic, so a least-degree monic annihilator exists. If $p$ and $q$ are two of them, then $(p-q)(T)=0$ and $\deg(p-q)<\deg p$.\\
If $p\neq q$, scaling $p-q$ gives a smaller monic annihilator, impossible. Hence $p=q$.\\
\hlt{Structural Commentary}: Existence comes from finite-dimensional dependence; uniqueness from minimality.\\
\hlt{Key Insight}: Two least monic annihilators cannot differ.
\end{proof}
\end{theorem}

\begin{definition} \label{def:operator-minimal-polynomial} \hlt{Minimal Polynomial of an Operator}\\
Let $\mathcal{V}\neq\{\mathbf{0}\}$ be finite-dimensional and let $T\in\LL(\mathcal{V})$. The \hlt{minimal polynomial} of $T$, denoted $\mu_T$, is the unique monic polynomial of smallest degree such that $\mu_T(T)=0$.
\end{definition}

\begin{definition} \label{def:minimal-polynomial-matrix} \hlt{Minimal Polynomial of a Matrix}\\
Let $A\in\operatorname{Mat}(n,n,\mathbb{F})$. The \hlt{minimal polynomial} of $A$, denoted $\mu_A$, is the unique monic polynomial of smallest positive degree such that $\mu_A(A)=0$.
\end{definition}

